% Source: public/notes/物工/問題/2023ap_exam.pdf, PDF pp. 8--9.
\documentclass[11pt,a4paper]{article}
\usepackage[margin=24mm]{geometry}
\usepackage{fontspec}
\usepackage{xeCJK}
\setCJKmainfont{HaranoAjiMincho-Regular.otf}
\setCJKsansfont{HaranoAjiGothic-Medium.otf}
\usepackage{amsmath,amssymb}
\setlength{\parindent}{1em}
\setlength{\parskip}{0.35em}
\title{物理工学専攻 院試解答集\\2023年 物理学 第4問}
\author{政岡凜太郎}
\date{}
\begin{document}
\maketitle

\section*{[1.1] 観測量の交換関係}
式(3)の定義と $[\hat a,\hat a^\dagger]=1$ を使う。$[\hat a^\dagger,\hat a]=-1$ なので
\[
[\hat x,\hat p]
=\frac{i}{2}[\hat a^\dagger+\hat a,\hat a^\dagger-\hat a]
=\frac{i}{2}\{-[\hat a^\dagger,\hat a]+[\hat a,\hat a^\dagger]\}
=\boxed{i}.
\]
$\hat x,\hat p$ は無次元の共役な観測量である。

\section*{[1.2] 真空のゆらぎ}
$\hat a|0\rangle=0$ と交換関係から、$\langle0|\hat a^\dagger\hat a|0\rangle=0$、$\langle0|\hat a\hat a^\dagger|0\rangle=1$ である。$\hat a^2$ と $(\hat a^\dagger)^2$ の真空期待値もゼロであり、$\langle\hat x\rangle=\langle\hat p\rangle=0$ だから
\[
\sigma_x^2=\langle\hat x^2\rangle=\frac12,
\qquad
\sigma_p^2=\langle\hat p^2\rangle=\frac12.
\]
従って $\boxed{\sigma_x\sigma_p=1/2}$。右辺も $\frac12|\langle[\hat x,\hat p]\rangle|=1/2$ なので、不確定性関係の等号が成立する。

\section*{[2.1] スクイーズ変換後の係数}
式(5)で $\hat S^\dagger\hat a\hat S$ と $\hat S^\dagger\hat a^\dagger\hat S$ を式(1)へ代入する。$\hat a\hat a^\dagger=\hat a^\dagger\hat a+1$ に注意すると
\begin{align*}
\hat S^\dagger\hat a^\dagger\hat a\hat S
&=\cosh(2r)\hat a^\dagger\hat a+\sinh^2r
-\frac{\sinh(2r)}2\bigl((\hat a^\dagger)^2+\hat a^2\bigr),\\
\hat S^\dagger\bigl((\hat a^\dagger)^2+\hat a^2\bigr)\hat S
&=\cosh(2r)\bigl((\hat a^\dagger)^2+\hat a^2\bigr)
-2\sinh(2r)\hat a^\dagger\hat a-\sinh(2r).
\end{align*}
よって式(6)の各係数は
\begin{align*}
\boxed{\omega(r)}&=\boxed{\omega\cosh(2r)-\Delta\sinh(2r)},\\
\boxed{\Delta(r)}&=\boxed{\Delta\cosh(2r)-\omega\sinh(2r)},\\
\boxed{E(r)}&=\boxed{\hbar\omega\sinh^2r-\frac{\hbar\Delta}{2}\sinh(2r)}.
\end{align*}
$r=0$ を入れると元の $\omega,\Delta,E=0$ に戻る。

\section*{[2.2] 基底状態と偶数光子}
$\omega>\Delta>0$ のとき $\Delta(r)=0$ とする実数 $r$ は
\[
\tanh(2r)=\frac{\Delta}{\omega},\qquad
\boxed{r=\frac14\log\frac{\omega+\Delta}{\omega-\Delta}}>0.
\]
この値で $\Omega=\sqrt{\omega^2-\Delta^2}>0$ とおけば $\omega(r)=\Omega$、$E(r)=\hbar(\Omega-\omega)/2$ であり、
\[
\hat S^\dagger(r)\hat H\hat S(r)
=\hbar\Omega\hat a^\dagger\hat a+\frac{\hbar}{2}(\Omega-\omega).
\]
右辺では $\hat a^\dagger\hat a\geq0$ なので $|0\rangle$ が基底状態である。元のハミルトニアンの基底状態とそのエネルギーは、それぞれ $\boxed{\hat S(r)|0\rangle}$、$\hbar(\Omega-\omega)/2$ となる。

$\hat S(r)$ の指数の各項 $\hat a^2,(\hat a^\dagger)^2$ は光子数を2だけ変える。真空は光子数0から始まるため、指数を展開して作用させても偶数光子の部分空間から出ない。従って $\hat S(r)|0\rangle$ には奇数光子成分がない。

\section*{[2.3] スクイーズされた基底状態のゆらぎ}
式(5)と観測量の定義から
\[
\hat S^\dagger\hat x\hat S=e^{-r}\hat x,
\qquad \hat S^\dagger\hat p\hat S=e^{r}\hat p.
\]
$\hat S|0\rangle$ での期待値を真空期待値に移すと、平均値はどちらもゼロであり、[1.2]の分散から
\[
\boxed{\sigma_x^2=\frac{e^{-2r}}2
=\frac12\sqrt{\frac{\omega-\Delta}{\omega+\Delta}}},
\qquad
\boxed{\sigma_p^2=\frac{e^{2r}}2
=\frac12\sqrt{\frac{\omega+\Delta}{\omega-\Delta}}}.
\]
$\Delta>0$ では真空に比べて $x$ の分散が小さく、$p$ の分散が大きい。一方、積は $\sigma_x\sigma_p=1/2=\frac12|\langle[\hat x,\hat p]\rangle|$ であり、等号は保たれる。$\Delta\to0$ では両分散は $1/2$ に戻る。

\section*{[3.1] ハイゼンベルク方程式の行列}
交換関係を積に適用すると
\[
[\hat a^\dagger\hat a,\hat a]=-\hat a,
\quad [(\hat a^\dagger)^2,\hat a]=-2\hat a^\dagger,
\quad [\hat a^2,\hat a^\dagger]=2\hat a.
\]
式(7)から $\dot{\hat a}=-i(\omega\hat a+\Delta\hat a^\dagger)$、$\dot{\hat a}^\dagger=i(\Delta\hat a+\omega\hat a^\dagger)$ となる。従って式(8)の行列は
\[
\boxed{M=\begin{pmatrix}\omega&\Delta\\-\Delta&-\omega\end{pmatrix}}.
\]

\section*{[3.2] 固有値と安定性}
$M^2=(\omega^2-\Delta^2)I$ なので固有値は
\[
\boxed{\lambda_\pm=\pm\sqrt{\omega^2-\Delta^2}}.
\]
$\omega>\Delta$ なら実数の $\lambda_\pm=\pm\Omega$ である。$M^2=\Omega^2I$を指数級数の偶数次と奇数次に適用すると、$e^{-iMt}=I\cos(\Omega t)-i(M/\Omega)\sin(\Omega t)$となり振動的かつ有界である。$\omega<\Delta$ なら$\gamma=\sqrt{\Delta^2-\omega^2}>0$として$M^2=-\gamma^2I$なので、同じ展開から$e^{-iMt}=I\cosh(\gamma t)-i(M/\gamma)\sinh(\gamma t)$となる。一方のモードは指数的に増大するので不安定である。境界 $\omega=\Delta$ では固有値がともにゼロでも $M^2=0$ かつ $M\ne0$ なので $e^{-iMt}=I-iMt$ は時間に比例して増大する。この意味でも有界な安定領域には含まれない。

\section*{[3.3] 境界での光子数}
$\omega=\Delta$ で $M^2=0$ を使い、初期条件 $\hat a(0)=\hat a$、$\hat a^\dagger(0)=\hat a^\dagger$ を式(8)に適用する。
\[
\begin{pmatrix}\hat a(t)\\\hat a^\dagger(t)\end{pmatrix}
=(I-iMt)\begin{pmatrix}\hat a\\\hat a^\dagger\end{pmatrix},
\qquad
\hat a(t)=(1-i\Delta t)\hat a-i\Delta t\hat a^\dagger.
\]
初期真空では $\langle\hat a^\dagger\hat a\rangle=0$、$\langle\hat a\hat a^\dagger\rangle=1$ である。従って $\hat a(t)$ の生成演算子の係数だけが $\langle0|\hat a^\dagger(t)\hat a(t)|0\rangle$ に寄与し、
\[
\boxed{\langle\hat n(t)\rangle=|{-i\Delta t}|^2=\Delta^2t^2}.
\]
$t=0$ でゼロに戻り、$\Delta t$ は無次元である。
\end{document}
